\chapter{Moderators in Temporal Networks}
\section{Time Series Models of Psychological Processes}
In order to study psychological processes as they unfold over time, researchers often use repeated sampling techniques, including longitudinal research designs, experience sampling methodology, and ecological momentary assessment. Latent variable models, such as growth curve models and cross-lagged panel models, are often used when few time points are available, but other techniques become available with a greater number of measurements. Specifically, \textit{time series analysis} can be applied when a large number of time points are available, and may afford a richer picture of the temporal dynamics governing intraindividual processes.

\subsection{Stationarity Assumptions}
Strict stationarity requires that all moments of a time series are independent of time---that is, they do not change over time. Second-order stationarity, however, only requires that the first and second order moments are independent of time. For normally-distributed (i.e., Gaussian) time series, all moments beyond the second order moments are fixed, and so second-order stationarity implies strict stationarity. If this is true, then it can be shown that: $\Ev{\by_t} = \bmu_t = \bmu$, which states that the mean does not change over time. Additionally, it must be true that:
\begin{equation}
    \bSigma(\by_t,\,\by_{t-k}) = \bSigma_k = \Ev{(\by_t - \bmu)(\by_{t-k} - \bmu)\tran},
\end{equation}
that is, the auto- and cross-covariances do not depend on occasion $t$ but depend on lag $k$ only, which is the fixed distance between occasions. 

\subsection{Estimating Temporal and Contemporaneous Networks}
In general, temporal networks are directed graphs with weighted edges that signify the magnitude and direction of time-lagged relationships between any pair of nodes. These visualizations essentially summarize the results of a multivariate VAR, wherein each node is regressed onto all other nodes at some time lag, and the resulting beta coefficients from each univariate model define the directed edges along with their weighted values. The contemporaneous network is generally modeled in conjunction with a temporal network, as it characterizes the inverse covariance matrix of the residuals across the different equations. 

\subsubsection{Graphical VAR}
This is the dominant method being used in the network psychopathology literature. First, imagine that you have $T$ repeated measures of $p$ variables for one subject, $n = 1$. The goal is to create a temporal network, which summarizes the time-lagged effects of each variable on every other variable, and then a contemporaneous network which is a GGM consisting of the partial correlations among the residuals of each time-lagged equation. The primary aim is to estimate these two networks together, such that the temporal network shows the lagged effects conditioned on the correlated residuals, and the contemporaneous network shows the relationships that exist once time-lagged effects have been controlled for. Furthermore, the idea behind the GVAR is to apply the LASSO regularization at every step of the way.

In the GVAR procedure, \cite{rothman2010sparse} propose an approach that they call "multivariate regression with covariance estimation' (MRCE). The algorithm starts by fitting intercept-only models in each of the $p$ equations, and then from those estimates (with zeros in place as coefficients for the predictor variables) finding the variance-covariance matrix of the residuals $\hat{\bm{\Sigma}}$. Then, they use penalized maximum likelihood to estimate the coefficient matrix, and the GLASSO to estimate the inverse residual covariances. This goes back and forth until the beta matrix converges to a pre-defined limit. The objective function being minimized is defined as:
\begin{equation}
    (\hat{\matr{B}}, \hat{\bTheta}) = \argmin_{\matr{B},\bTheta} \lb g(\matr{B},\bTheta) + \lambda_1 \sum_{j' \neq j}|\theta_{j'j}| + \lambda_2 \sum_{j = 1}^p \sum_{k = 1}^q |b_{jk}| \rb.
\end{equation}
This essentially follows the process of (1) Starting with a fixed B or Theta (the former may be estimated by separate lasso regressions; the latter may be the error covariance), then estimate either Theta or B (the opposite of what you started with), and repeat until convergence. 

%\apalevelfour{Partial contemporaneous correlations}
The graphical VAR model estimates $\matr{B}$ and $\bTheta$ using a penalized likelihood approach, wherein $\lambda_1$ and $\lambda_2$ are employed as tuning parameters for penalizing the estimates within each matrix. This method uses a criterion-based model selection procedure for choosing optimal values $\lambda_1$ and $\lambda_2$. Then, the strengths of the links in the graphical model are assessed by computing so-called \textit{partial directed correlations} (PDC) and \textit{partial contemporaneous correlations} (PCC) as measures of strength of the lagged and same-time associations between variables, respectively \cite{wild2010graphical}. 

The PCC is defined as the correlation between two variables at the same point in time after removing the linear effects of the other variables at the same point in time and all variables at previous time points. These values are computed based on elements of the estimated residual covariance matrix $\hat{\matr{\Theta}}$. This is the same procedure as constructing a GGM with cross-sectional data.
\begin{equation}
    \text{PCC}(X_{i,t}, X_{j,t}) = - \frac{\hat{\theta}_{ij}}{\sqrt{\hat{\theta}_{ii} \hat{\theta}_{jj}}}.
\end{equation}

%\apalevelfour{Partial directed correlations}
Since the regression coefficients in $\matr{B}$ depend on the unit of measurement, the strength of lagged associations among variables can be measured as PDCs. These values measure the linear association between a dependent variable at time $t$ and an explanatory variable at time $t - 1$ after removing the linear effects of all other variables at time $t-1$. They therefore quantify the direct influence of the explanatory variable on the dependent variable. The PDC is calculated by rescaling the autoregressive coefficients $\hat{\beta}_{ij} \in \matr{B}$
\begin{equation}
    \text{PDC}(X_{i,t}, X_{j, t-1}) = \frac{\beta_{ij}}{\sqrt{\Sigma_{ii} \Theta_{jj} + \beta_{ij}^2}}.
\end{equation}

\subsection{Seemingly Unrelated Regression (SUR)}
SUR models have been widely used within econometrics, and as their name implies: This method is used when the goal is to perform multivariate regression, wherein each equation appears independent from the others, but really the residuals are correlated across them all. The general approach to estimating these kinds of networks has been to use an iterative procedure that relies on \textit{generalized least squares (GLS)} rather than OLS. The reason for this is because of the correlated residuals across equations. This approach is frequently utilized within the `seemingly unrelated regression' (SUR) framework, and that forms the basic of the graphical VAR (GVAR) as it has been used within the psychological network literature.

\subsubsection{Coping with correlated errors}
Take the standard case of linear regression $\matr{y} = \matr{X}\beta + \epsilon$, where we obtain an unbiased estimate $\hat{\beta}$ of the variable coefficients by solving: $\hat{\beta} = (\matr{X}\tran \matr{X})^{-1} \matr{X}\tran \matr{y}$. This is known as the ordinary least squares (OLS) estimate of $\beta$, but is only efficient if there are no correlations among the residuals. One way of dealing with this challenge is by using \textit{generalized least squares (GLS)} to create $\hat{\beta}$, where the inverse of the error covariance matrix is used: $\hat{\beta}_{GLS} = (\matr{X}\tran \hat{\matr{\Sigma}}^{-1} \matr{X})^{-1} \matr{X}\tran \hat{\matr{\Sigma}}^{-1} \matr{y}$, where the inverse covariance error matrix is $\hat{\matr{\Sigma}}^{-1}$.

\subsubsection{Two-stage graphical SUR}
In general, my proposal for fitting this type of model focuses on (a) utilizing the SUR framework to fit the models of interest, and (b) separating the variable-selection procedure from the final estimation in order to hopefully improve inference and generalizability. 

%\apalevelfour{Step 1: Variable selection} 
To do this, use the hierarchical LASSO as per the \code{glinternet} (or \code{hierNet}) package in the R language. Do this in sequential, nodewise fashion until all equations are complete.

%\apalevelfour{Step 2: GLS model fitting}
Next, refit the constrained system of equations (with restrictions determined in Step 1) using standard GLS for multivariate regression.

\begin{enumerate}
    \item Perform variable selection before model fitting
    \item Do this in sequential, nodewise fashion.
    \item The selected variables for each equation are then used to construct the SUR model. 
    \item Advantages/similarities to GVAR, beyond the added features of including moderators.
    \item Could be extended to multilevel modeling.
\end{enumerate}

\section{Experience Sampling Methodology}
To study moment-to-moment changes in psychological and affective experience, researchers have utilized \textit{intensive longitudinal modeling} to better understand the dynamics of daily life. Some methods have included computing cross-correlations, average weighted coherence indices using cross-spectral analysis \cite{sadler2009we}, latent growth curve models (LGCMs) and multilevel models (MLMs) with time-varying covariates \cite{dermody2017modeling}. However, these tend to involve aggregating across longitudinal measurements to summarize variability; less frequently do these methods actually characterize the \textit{patterns} of temporal variation, and \textit{changes} in associations over time. 

LGCMs are designed to characterize individual differences in growth trajectories over time \cite{bollen2006lgcms}. \cite{wright2013parallel} showed that individuals linearly increase in avoidant personality features and neuroticism over time, while decreasing in dominance and affiliation. In essence, these methods (e.g., LGCMs) can be used to demonstrate associated change, but not changes in association \cite{dermody2017modeling}. Polynomial models can capture a bit more detail in this type of change (e.g.) \cite{wright2015daily}, but still fail to capture the dynamic relationships among variables. Even MLMs with time-varying predictors produce coefficients estimating average moment-to-moment associations. 

\subsection{Time-Varying Effect Models}
For the above reasons, some researchers have suggested a shift towards focusing on what's known as a \textit{time-varying effect model} (TVEM, e.g.) \cite{dermody2017modeling,tan2012time}. TVEM is an extension of MLM that uses coefficient functions to describe dynamic associations between variables via the estimation of non-parametric splines---e.g., polynomials defined over specific intervals of time. A powerful capability that these models have is to specify when moderating influences have the greatest effect on other variables of interest e.g. \cite{wright2014treating,shiyko2012using}.

\cite{dermody2017modeling} used this method to study relationship partners having discussions in real time, where video records were coded using joysticks to operationalize interpersonal dynamics (e.g., expressions of affiliation and control) at a rate of 2 Hz. The time-varying associations between the wives' and husbands' interpersonal behaviors were captured using two separate models

\section{Multivariate Time Series Models}
Time series analysis spans a wide range of different modeling approaches. In general, the goal in univariate time series analysis is to study how a particular process changes over time, as well as quantify the extent to which its behavior depends on characteristics of its previous states. The most basic model of this kind is the \textit{autoregressive (AR) model}, wherein the value of a given variable is regressed on measurements of its past states. The multivariate extension of this model is known as a \textit{vector autoregressive model (VAR)}, and is often the starting point for constructing temporal networks from intensive longitudinal data.


\subsection{Autoregressive Models}
Lag 1 model AR(1)
\begin{equation}
    X_t = \beta_0 + \beta_1 X_{t - 1} + \epsilon_t
\end{equation}
Model with $p$ lags
\begin{equation}
    X_t = \beta_0 + \beta_1 X_{t - 1} + \cdots + \beta_p X_{t - p} + \epsilon_t
\end{equation}

\subsection{Vector Autoregressive Model}
Ordinary least squares
\begin{equation}
    \matr{\hat{\beta}} = (\matr{X}\tran \matr{X})^{-1} \matr{X}\tran \matr{y}
\end{equation}

\subsection{Graphical VAR}

\subsection{Seemingly Unrelated Regressions}
Generalized least squares
\begin{equation}
    \begin{split}
        \by &= \bX\bbeta + \epsilon \\
        \hat{\bbeta} &= (\bX\tran \bSigma\inv \bX)\inv \bX\tran \bSigma\inv \by
    \end{split}
\end{equation}

Here I'll talk about SUR models, and how they are useful for estimating VARs, including mmVARs. But first, I will display a similarity measure that I'm using to compare the \textit{true} (i.e., pre-specified) and \textit{estimated} model matrices that are of interest within this framework. Namely, we are interested in estimating the \textit{beta matrix} $\bm{\beta}$, and the \textit{residual covariance matrix} $\bm{\Sigma}$. Here is one metric that can be used to quantify the degree of similarity between two matrices (e.g., the true matrix $\bm{\beta}$ and the matrix $\hat{\bm{\beta}}$ that is estimated by the SUR model). The metric is the same whether it's written in matrix- or vector-form.
\begin{equation}
    \text{cosine similarity} = \cos({\matr{A},\matr{B}}) = \frac{\tr(\matr{A}\tran \matr{B})}{\norm{\matr{A}}_{\text{F}} \norm{\matr{B}}_{\text{F}}} = \frac{\displaystyle \sum_{i = 1}^n a_i b_i}{\sqrt{\displaystyle \sum_{i = 1}^n a_i^2} \sqrt{\displaystyle \sum_{i = 1}^n b_i^2}}
\end{equation}
where the \textit{Frobenius norm} is equal to: $\norm{\matr{A}}_{\text{F}} = \sqrt{\tr(\matr{A}\tran\matr{A})} = \sqrt{\displaystyle \sum_{i = 1}^m \sum_{j = 1}^n |a_{i,j}|^2}$.

\subsection{Two-Stage Graphical SUR}
Something something

\section{Mixed Moderated Vector Autoregression}
First, we need to consider how we will be conceptualizing vector autoregressive (VAR) models. The most basic option is to specificy a VAR with no additional characteristics, which forces us to adopt the following assumptions: (a) the error terms are uncorrelated over time, and (b) it's getting dark in this tomato. Here is the almighty VAR general equation:
\begin{equation}
    \matr{y}_t = \matr{A}_1 \matr{y}_{t-1} + \cdots + \matr{A}_p \matr{y}_{t-p} + \matr{u}_t
\end{equation}
where $A_i$ is a $(K \times K)$ coefficient matrix, and $u_t = (u_{1t},\ldots,u_{Kt})\tran$ is an unobservable error term. $u_t$ is assumed to be a zero-mean independent white noise process with time-invariant, positive definite covariance matrix $E(u_t,u_t\tran) = \Sigma_u$. I.e., $u_t$ is an independent stochastic vector with $u_t \sim (0, \Sigma)$. The process is \textit{stable} if
\begin{equation}
    \det(I_K - A_1 z - \cdots - A_p z^p) \neq 0 \;\; \text{for} \;\; |z| \leq 1
\end{equation}
We could also add a linear trend, or any $q$-order polynomial term if we wish to detrend the series and condition the autoregressive and cross-lagged parameters on those effects, for instance:
\begin{equation}
    \matr{y}_t = v_0 + v_1 t + \matr{A}_1 \matr{y}_{t-1} + \cdots + \matr{A}_p \matr{y}_{t-p} + \matr{u}_t
\end{equation}

We also have the \textit{structural VAR}, which is estimated in terms of a generalized VECM
\begin{equation}
    \matr{A}_0 \Delta \matr{y}_t = \bPi^* \by_{t-1} + \bGamma_1^* \Delta \by_{t-1} + \cdots + \bGamma_{p-1}^* \Delta \by_{t - p + 1} + \matr{C}^* \matr{D}_t + \matr{B}^* z_t + v_t
\end{equation}
while the following shows the SVAR in terms of a standard VAR model
\begin{equation}
    \matr{A}_0 \matr{y}_t = \matr{A}_1^* \matr{y}_{t-1} + \cdots + \matr{A}_p^* \matr{y}_{t - p} + \matr{B} \epsilon_t
\end{equation}
with the reduced form
\begin{equation}
    \matr{y}_t = \matr{A}_1 \matr{y}_{t-1} + \cdots + \matr{A}_p \matr{y}_{t-p} + \matr{u}_t, \;\; \text{where} \;\; \matr{u}_t = \matr{A}_0^{-1} \matr{B} \epsilon_t, \;\; \text{and} \;\; \matr{A}_0 = \matr{I}_K
\end{equation}

I'll show the VAR model here, but our first task is to consider how we will estimate this model. In simple cases, the model parameters can be solved for via maximum likelihood estimation, where the following defines this procedure in terms of matrix algebra. Using this method entails estimating the autoregressive matrix as a whole, whereas an alternative method might be to conduct separate regressions and aggregate the results to form a weighted adjacency matrix. To my knowledge, the former method is not possible in cases that involve mixed conditional distributions. In such instances, \textit{nodewise regression} provides a useful alternative and affords the modeling of complex data structures. 

Once method of estimating edges is determined, an important decision for interpreting the graph is whether or not a sparse solution is desired. That is, we may wish to avoid confirmatory interpretations of the results until we have employed some decision rule for identifying important relationships among the variables of interest. Next, here are a bunch of regularization and decision-rule techniques that are often used within psychological applications of graph-theoretic models. For instance, here is penalized likelihood of the precision matrix for covariance-estimation methods:
\begin{equation}
    l(\bTheta) = \log \det \bTheta - \tr(\matr{S}\bTheta) - \lambda_p \sum_{i \neq j} |\bTheta_{i,j}|
\end{equation}

\subsection{Maximum Likelihood Approaches}
Let $X_n$ be the observed data for the \textit{i}th $(i \in 1,\ldots,p)$ variable at the \textit{n}th $(n \in 1,\ldots,N)$ observation. First, we compute the $p \times p$ covariance matrix defined by the MLE as
\begin{equation}
    \bSigma = \frac{1}{N} \sum_{n=1}^N (X_n - \Bar{X})(X_n - \Bar{X})\tran, \; \Bar{X} = \frac{1}{N} \sum_{n=1}^N X_n
\end{equation}
Then, we can use $\bSigma$ to estimate the precision matrix $\bTheta$ as: $\hat{\bSigma}\inv = \hat{\bTheta}$. 

The precision matrix $\hat{\bTheta}$ then contains the covariances $\hat{\theta}_{ij}$ and variances $\hat{\theta}_{ii}$.
\begin{equation}
    \hat{\bTheta} = 
    \begin{bmatrix}
    \hat{\theta}_{ii} &  &  \\
    \vdots & \ddots &  \\
    \hat{\theta}_{ij} & \cdots & \hat{\theta}_{jj}
    \end{bmatrix}
\end{equation}
The partial correlations are then obtained as
\begin{equation}
    \hat{\rho}_{ij} = \frac{-\hat{\theta}_{ij}}{\sqrt{\hat{\theta}_{ii}\hat{\theta}_{jj}}}
\end{equation}

Now, if we wish to reduce the prevalence of spurious relationships, we must set a decision rule for determining whether $\hat{\rho}_{ij}$ should be fixed to 0. To do this in accordance with standard hypothesis testing procedures, we can use Fisher's \textit{Z}-transformation:
\begin{equation}
    z_{ij} = \frac{1}{2} \log \Big(\frac{1 + \hat{\rho}_{ij}}{1 - \hat{\rho}_{ij}} \Big)
\end{equation}
The resultant values can now be interpreted with reference to the standard \textit{z} distribution, which is approximated as
\begin{equation}
    z_{ij} \sim \mathcal{N} \Bigg(\frac{1}{2} \log \Big(\frac{1 + \hat{\rho}_{ij}}{1 - \hat{\rho}_{ij}} \Big), \: \frac{1}{n - 3 - s} \Bigg)
\end{equation}
where $s$ denotes the number of variables controlled for $(p - 1)$, and the standard error is defined as $\sqrt{\frac{1}{n - 3 - s}}$.

Given our \textit{hard-threshold} choice of the false-positive rate $\alpha$, we can define confidence intervals based on the corresponding critical values $Z_{\alpha/2}$. The confidence interval for $z_{ij}$ can then be defined as:
\begin{equation}
    \begin{split}
        Z_{Lower} &= z_{ij} - Z_{\alpha/2} \sqrt{\frac{1}{n - 3 - s}} \\
        Z_{Upper} &= z_{ij} + Z_{\alpha/2} \sqrt{\frac{1}{n - 3 - s}}
    \end{split}
\end{equation}
Finally, to obtain the interval for $\hat{\rho}_{ij}$ we must transform back to the relevant scale
\begin{equation}
    \hat{\rho}_{ij_{L}} = \frac{\exp(2 Z_{L}) - 1}{\exp(2 Z_{L}) + 1} \;\; \text{and} \;\; \hat{\rho}_{ij_{U}} = \frac{\exp(2 Z_{U}) - 1}{\exp(2 Z_{U}) + 1}
\end{equation}

\subsection{Partial Contemporaneous and Directed Correlations}
The graphical VAR model estimates $\matr{B}$ and $\bTheta$ using a penalized likelihood approach, wherein $\lambda_1$ and $\lambda_2$ are employed as hyperparameters for penalizing the estimates within each matrix. This appears to be framed as an alternative to the SVAR model, and utilizes a criterion-based model selection procedure for choosing optimal values $\lambda_1$ and $\lambda_2$. Then, the strengths of the links in the graphical model are assessed by computing so-called \textit{partial directed correlations} (PDC) and \textit{partial contemporaneous correlations} (PCC) as measures of strength of the lagged and same-time associations between variables, respectively \cite{wild2010graphical}. 

The PCC is defined as the correlation between two variables at the same point in time after removing the linear effects of the other variables at the same point in time and all variables at previous time points. These values are computed based on elements of the estimated residual covariance matrix $\hat{\bTheta}$. These are literally just partial correlations...
\begin{equation}
    \text{PCC}(X_{i,t}, X_{j,t}) = - \frac{\hat{\theta}_{ij}}{\sqrt{\hat{\theta}_{ii} \hat{\theta}_{jj}}}
\end{equation}
Since the regression coefficients in $\matr{B}$ depend on the unit of measurement, the strength of lagged associations among variables can be measured as PDCs. These values measure the linear association between a dependent variable at time $t$ and an explanatory variable at time $t - 1$ after removing the linear effects of all other variables at time $t-1$. They therefore quantify the direct influence of the explanatory variable on the dependent variable. The PDC is calculated by rescaling the autoregressive coefficients $\hat{\beta}_{ij} \in \matr{B}$
\begin{equation}
    \text{PDC}(X_{i,t}, X_{j, t-1}) = \frac{\beta_{ij}}{\sqrt{\Sigma_{ii} \Theta_{jj} + \beta_{ij}^2}}
\end{equation}